% Part: many-valued-logic
% Chapter: three-valued-logics
% Section: lukasiewicz

\documentclass[../../../include/open-logic-section]{subfiles}

\begin{document}

\olfileid[pt]{mvl}{inf}{god}

\olsection{Lógicas de G\"odel}

\begin{editorial}
  Este é um breve ``esboço'' de uma seção sobre a lógica de infinitos
  valores de G\"odel.
\end{editorial}

\begin{defn}\ollabel{def:goedel} A lógica de infinitos valores de
G\"odel~$\LogGod[\infty]$ é definida pela matriz:
\begin{enumerate}
  \item A linguagem proposicional padrão $\Lang L_0$ com
  $\lfalse$, $\lnot$, $\land$, $\lor$, $\lif$.
  \item O conjunto de valores de verdade $V_\infty$.
  \item $1$ é o único valor designado, O.s.\@, $V^+ = \{1\}$.
  \item As funções de verdade são dadas pelas seguintes funções:
  \begin{align*}
    \tf{\lfalse} & = 0\\
    \tf{\lnot}[\LogGod](x) & = \begin{cases}
      $1$ & \text{se } x =0\\
      $0$ & \text{caso contrário}
    \end{cases}\\
    \tf{\land}[\LogGod](x,y) & = \min(x,y)\\
    \tf{\lor}[\LogGod](x,y) & = \max(x,y)\\
    \tf{\lif}[\LogGod](x,y) & = \begin{cases}
      1 & \text{se } x \le y\\
      y & \text{caso contrário.}
    \end{cases}
    \end{align*}
\end{enumerate}
A lógica $m$-valorada de G\"odel é definida da mesma maneira, exceto
que $V = V_m$.
\end{defn}

\begin{prop}
  A lógica $\LogGod[3]$ definida por \olref[thr][god]{defn:goedel}
  é a mesma que $\LogGod[3]$ definida por \olref{def:goedel}.
\end{prop}

\begin{proof}
  Isso pode ser visto comparando as tabelas de verdade dos conectivos
  dadas em \olref[thr][god]{defn:goedel} com as tabelas de verdade
  determinadas pelas equações em \olref{def:goedel}:
  \begin{center}
    \begin{tabular}{c|c} 
      $\tf{\lnot}[\LogGod[3]]$ & \\ 
      \hline  
      $1$ & $0$ \\ 
      $1/2$ & $0$ \\
      $0$ & $1$ 
    \end{tabular}
    \quad
    \begin{tabular}{c|ccc} 
      $\tf{\land}[\LogGod]$ & $1$ & $1/2$ & $0$ \\ 
      \hline 
      $1$ & $1$ & $1/2$ & $0$ \\ 
      $1/2$ & $1/2$ & $1/2$ & $0$\\ 
      $0$ & $0$ & $0$ & $0$ 
    \end{tabular}
    \\[2ex]
    \begin{tabular}{c|ccc} 
      $\tf{\lor}[\LogGod]$ & $1$ & $1/2$ & $0$ \\ 
      \hline 
      $1$ & $1$ & $1$ & $1$ \\ 
      $1/2$ & $1$ & $1/2$ & $1/2$ \\
      $0$ & $1$ & $1/2$ & $0$ 
    \end{tabular}
    \quad
    \begin{tabular}{c|ccc} 
      $\tf{\lif}[\LogGod]$ & $1$ & $1/2$ & $0$ \\ 
      \hline 
      $1$ & $1$ & $1/2$ & $0$ \\ 
      $1/2$ & $1$ & $1$ & $0$  \\ 
      $0$ & $1$ & $1$ & $1$ 
    \end{tabular}
  \end{center} 
\end{proof}

\begin{prop}\ollabel{prop:god-infty-m}
  Se $\Gamma \Entails[\LogGod[\infty]] !B$, então $\Gamma
  \Entails[\LogGod[m]] !B$ para todo~$m \ge 2$.
\end{prop}

\begin{proof}
  Exercício.
\end{proof}

\begin{prob}
  Demonstre \olref[mvl][inf][god]{prop:god-infty-m}.
\end{prob}

Na verdade, a recíproca também vale.

Assim como $\LogGod[3]$, $\LogGod[\infty]$ tem como tautologias todas
as !!{formula}s válidas intuicionisticamente, e os mesmos exemplos de
não tautologias são não tautologias de~$\LogGod[\infty]$:
\begin{align*}
  & p \lor \lnot p && (p \lif q) \lif (\lnot p \lor q) \\
  & \lnot\lnot p \lif p && \lnot(\lnot p \land \lnot q) \lif (p \lor q) \\
  & ((p \lif q) \lif p) \lif p && \lnot(p \lif q) \lif (p \land \lnot q)
\end{align*}
O exemplo de uma !!{formula} inválida intuicionisticamente que, não
obstante, é uma tautologia de~$\LogGod[3]$,
$(p \lif q) \lor (q \lif p)$, também é uma tautologia em
$\LogGod[\infty]$. Na verdade, $\LogGod[\infty]$ pode ser
caracterizada como a lógica intuicionista à qual se acrescenta o
esquema $(!A \lif !B) \lor (!B \lif !A)$. Isso foi demonstrado por
Michael Dummett e, por isso, $\LogGod[\infty]$ é frequentemente chamada
lógica de G\"odel--Dummett~$\Log{LC}$.

\begin{prob}
  Mostre que $(p \lif q) \lor (q \lif p)$ é uma
  tautologia de~$\LogGod[\infty]$.
\end{prob}

\begin{prob}
  Mostre que $(p \lif q) \lor (q \lif r) \lor (r \lif s)$, que é
  uma tautologia de $\LogGod[3]$, não é uma tautologia
  de~$\LogGod[\infty]$.
\end{prob}

\end{document}
